\scanpage{98}
\mspage{13}
% archive sheet number: 97

si $g$ est un générateur topologique de $I_t=I/P$ (\unsure{horreur!}), \inserted{$g_E$ son action sur $E$,} on a la relation
\[
(\mathrm{id}-g_E^{m'})^{i+1}=0
\tag*{(3.3)}
\]

\uline{Remarques} 3.4. a) La démonstration que nous allons donner \inserted{au \deleted{\unsure{4}}5} \unsure{marche|montre} également, mutatis mutandis, dans le cas analytique complexe, pour un morphisme propre $f:X\mathbin{-}S$ d'espaces analytiques, $S$ régulier de dim 1. \inserted{(auquel cas il n'y a plus lieu d'introduire \unsure{le groupe} $P$!)} L'énoncé dans ce cas est dû à P.A. Griffiths.
\margin{plus haut}
% slanted; bracket alongside paragraph a)

b) La démonstration qui suit \deleted{\unsure{montre que}} \inserted{donne en} fait une filtration \uline{canonique} de $E$, qui sera de plus invariante sous l'action de $\unsure{\pi|\Pi}$ tout entier, pas seulement de $I$. Elle sera étudiée plus bas à un point de vue arithmétique.

Signalons tout de suite :

\uline{Corollaire 3.5} \deleted{Supposons que} Soit \deleted{$\unsure{X}$} \inserted{$n\geq0$, et} \deleted{$U$ lisse quasi-projectif de type fini sur $K$\inserted{de dim \deleted{relative} $n$},} supposons \deleted{que}

(i) Pour tte extension finie $K'$ de $K$, si $S'$ est le normalisé de $S$ dans $K'$, la \uline{résolution des singularités forte} est possible pour \unsure{tt|$\forall$} schéma de type fini $X'$ sur $\hat S'$ de dim relative $\leq n$.

(ii) Pour tout \deleted{$\unsure{X'|K'},\unsure{S'}$} $X'$ comme dessus, avec $X'$ régulier, et tout diviseur régulier $Y'$ de $X'$, le th.
